\subsection{Geometric Properties of Springer fibers}

\subsubsection{Connectivity} The Springer fibers $\sB_{e}$ are connected. See \cite[p.131, Prop 1]{StConj}, \cite[Ch II, Cor 1.7]{Spaltenstein}, and Exercise \ref{exI:conn}.

\subsubsection{Equidimensionality} Spaltenstein \cite{Spa77}, \cite[Ch II, Prop 1.12]{Spaltenstein} showed that all  irreducible components of $\sB_{e}$ have the same dimension. 

\subsubsection{The dimension formula} Let $d_{e}$ be the dimension of $\sB_{e}$. Steinberg \cite[Thm 4.6]{St} and Springer \cite{Springer} showed that
\begin{equation}\label{Spr dim}
d_{e}=\frac{1}{2}(\dim\sN-\dim G\cdot e)=\frac{1}{2}(\dim C_{G}(e)-r).
\end{equation}


\subsubsection{Centralizer action} The group $G\times \Gm$ acts on $\sN$ (where $\Gm$ by dilation). Let $\tilG_{e}=\Stab_{G\times \Gm}(e)$ be the stabilizer. Then $\tilG_{e}$ acts on $\sB_{e}$. Note that $\tilG_{e}$ always surjective onto $\Gm$ with kernel $C_{G}(e)$.

The action of $\tilG_{e}$ on $\cohog{*}{\sB_{e}}$ factors through the finite group $A_{e}=\pi_{0}(\tilG_{e})=\pi_{0}(G_{e})$. Note that $A_{e}$ depends not only on the isogeny class of $G$, but on the isomorphism class of $G$. For example, for $e$ regular, $A_{e}=\pi_{0}(ZG)$ where $ZG$ is the center of $G$. The action of $A_{e}$ on $\cohog{*}{\sB_{e}}$ further factors through the image of $A_{e}$ in the adjoint group $G^{\ad}$.


\subsubsection{Purity} Springer \cite{SpringerPure} proved that the cohomology of $\sB_{e}$ is always pure (in the sense of Hodge theory when $k=\CC$, or in the sense of Frobenius weights when $k$ is a finite field). 


\subsubsection{} Let $e\in\sN$. Consider the restriction map $i^{*}_{e}: \cohog{*}{\sB}\to \cohog{*}{\sB_{e}}$ induced by the inclusion $\sB_{e}\incl\sB$. Then the image of $i^{*}_{e}$ is exactly the invariants of $\cohog{*}{\sB_{e}}$ under $A_{e}$. This is a theorem of Hotta and Springer \cite[Theorem 1.1]{HottaSpringer}. In particular, when $G$ is of type $A$, such restriction maps are always surjective.


\subsubsection{Parity vanishing} De Concini, Lusztig and Procesi \cite{DLP} proved that $\cohog{i}{\sB_{e}}$ vanishes for all odd $i$ and any $e\in\sN$. When $k=\CC$, they prove a stronger statement: $\cohog{i}{\sB_{e},\ZZ}$ vanishes for odd $i$ and is torsion-free for even $i$.



